\documentclass[10pt,a4paper]{amsart}
\usepackage[margin=19mm]{geometry}
\usepackage{amsmath,amssymb,mathtools}
\usepackage[scr=rsfs,cal=euler]{mathalfa}
\usepackage{xfrac}
\usepackage[unicode,hidelinks,bookmarks=false]{hyperref}
\newcommand{\ie}{i.e.\ }
\newcommand{\cf}{cf.\ }
\newcommand{\resp}{respectively\ }
\newcommand{\Subsection}{\subsection}
\providecommand{\nobreakdash}{\nobreak\hbox{-}}
\setlength{\emergencystretch}{2em}
\multlinegap0pt
\usepackage{bm}
\usepackage{bbm}
\usepackage{dsfont}
\usepackage{xspace}
\usepackage{cancel}
\usepackage{mathtools}
\usepackage{amsthm}
\makeatletter
\@ifundefined{dfn}{\newtheorem{dfn}{Dfn}}{}
\@ifundefined{cor}{\newtheorem{cor}[dfn]{Corollary}}{}
\@ifundefined{prop}{\newtheorem{prop}[dfn]{Proposition}}{}
\@ifundefined{thm}{\newtheorem{thm}[dfn]{Theorem}}{}
\makeatother
\title[Motivic Modularity of CM K3 Surfaces]{Motivic Modularity of CM K3 Surfaces}
\author{Rikuto Ito}
\date{Source: arXiv:2406.16325v4 (CC BY 4.0)}
\begin{document}
\maketitle

\noindent\emph{Source and licence.} Motivic Modularity of CM K3 Surfaces. Version arXiv:2406.16325v4 (CC BY 4.0), distributed under the Creative Commons Attribution 4.0 International licence, as recorded in the arXiv licence metadata for this exact version and shown on the abstract page https://arxiv.org/abs/2406.16325v4. Full rights notice: licenses/arxiv-2406.16325-CC-BY-4.0.txt.\par\smallskip

\noindent\textit{Editorial scope.} Counted unit: the Prop whose source label is \texttt{onegai} in the article (source0.tex lines 550--552, source label \texttt{onegai}), delivered together with the article's own complete original proof of it (source0.tex lines 553--570, one \texttt{proof} environment of 18 lines). No mathematics is written, repaired or re-derived: statement and proof are the article's own text line for line. Required context retained uncounted: main theorem (lines 82--93); hecke1 (lines 346--352); simple cm1 (lines 538--549). Every definition, display and label of those lines is retained unchanged. Omitted from this package and not redistributed: 1--81, 94--345, 353--537, 571--610. Nothing inside a retained excerpt is omitted or altered: every retained body is delivered exactly as the article's lines are written, so no omission receipt of any kind accompanies this package. Citations: the retained excerpts carry no citation command at all, so no bibliography entry of the article is transcribed and no reference is redistributed. Reference adaptation: none was needed and none was made; every reference inside the delivered unit resolves inside it, so no unresolved reference or citation is produced. Printed slips kept literal: no printed word, symbol, hypothesis or number of the counted statement or of its proof was altered. Layout and class: the article's own class is \texttt{amsart} with packages including fullpage, amsmath, amssymb, mathrsfs, xy, amsthm, graphicx; the wrapper is the lane's pinned \texttt{amsart} editorial wrapper at 10pt on A4, which restates the theorem-like environments the retained text needs and adds the article's own macro definitions that the retained text needs (none), with extra wrapper packages (bm, bbm, dsfont, xspace, cancel, mathtools). The delivered numbering is the unit's own. Assets: no retained span contains a figure environment, a graphics-inclusion command, a PDF-page inclusion, a table or a TikZ picture; no image file is copied into this package, the package declares no asset dependency and no image file is redistributed with it. Licence: version arXiv:2406.16325v4 of this article is distributed under the Creative Commons Attribution 4.0 International licence, as recorded in the arXiv licence metadata for this exact version and shown on the abstract page https://arxiv.org/abs/2406.16325v4. A full rights notice is delivered at licenses/arxiv-2406.16325-CC-BY-4.0.txt. Characters: every retained excerpt is pure ASCII, so the delivered engine, XeLaTeX with its default Latin Modern text font, reports no missing character.
    \begin{thm}\label{main theorem} 
        Let $X$ be a K3 surface defined over a number field $k\subset \mathbb{C}$. Suppose that $X$ has CM over $k$ by a  CM field $E$, i.e., $X_{\mathbb{C}}$ is of CM type and $k\supset E$. Then there exists a unique locally constant homomorphism $u : \mathbb{A}^{\times}_{k} \to E^{\times}\subset \mathbb{A}^{\times}_{E}$ satisfying the following properties (i)-(iii):
        \\ \indent (i) Fix an embedding $\tau : E\to \mathbb{C}$. The function \begin{align}
            \chi_{\tau} : \mathbb{A}^{\times}_{k}\xrightarrow{u(-)\cdot \overline{{\rm Nm}_{k/E}(-)}/{\rm Nm}_{k/E}(-)} \mathbb{A}^{\times}_{E}\xrightarrow{\tau-{\rm projection}} \mathbb{C}^{\times}\notag\end{align}
         is an algebraic Hecke character, where ${\rm Nm}_{k/E} : \mathbb{A}^{\times}_{k}\to \mathbb{A}^{\times}_{E}$ denotes the norm map for the extension $k/E$, and $\overline{(-)} : \mathbb{A}^{\times}_{E}\to \mathbb{A}^{\times}_{E}$ denotes the complex conjugation.
        \\\indent (ii) Let \begin{align}\rho_{l} : G_{k}\to GL(T_{l}(X_{\overline{k}})\otimes \mathbb{Q}_{l}) \notag
        \end{align}
    be the $(22-\rho(X_{\overline{k}}))$-dimensional $l$-adic representation. Then $\rho_{l}$ is diagonalized by the algebraic Hecke characters $\chi_{\tau}$, where $\tau$ runs over ${\rm Hom}(E, \mathbb{C})$.
    \\ \indent (iii) Fix an embedding $\tau\in {\rm Hom}(E, \mathbb{C})$ and a prime ideal $\mathfrak{p}\subset k$ with $(l, \mathfrak{p})=1$. Then the following conditions are equivalent: 
     \\(a) The $l$-adic representation $\rho_{l}$ is unramified at $\mathfrak{p}$, i.e., if $I_{\mathfrak{p}}$ denotes the inertia group of $\mathfrak{p}$, then $\rho_{l}(I_{\mathfrak{p}})=\{id\}$.
     \\(b) The Hecke character $\chi_{\tau}$ is unramified at $\mathfrak{p}$, i.e., if $\mathscr{O}_{\mathfrak{p}}$ is the ring of integers of $k_{\mathfrak{p}}$, then  $\chi_{\tau}(\mathscr{O}^{\times}_{\mathfrak{p}})=1.$
       \end{thm}

    \begin{prop}\label{hecke1}
        Suppose that $X$ is a K3 surface with CM over a number field $k\subset \mathbb{C}$. Let $E$ be the reflex field. Then there exists a unique locally constant homomorphism $ u: \mathbb{A}^{\times}_{k}\to E^{\times}$ satisfying the following condition: 
        \\ For all $t\in \mathbb{A}^{\times}_{k}$, 
        \begin{align}
            \rho^{ab}({\rm art}_k(t))=u(t)\overline{{\rm Nm}_{k/E}(t)}/{\rm Nm}_{k/E}(t)\in G(\mathbb{A}_{f})\notag.
        \end{align}
    \end{prop}

      \begin{cor}\label{simple cm1} Let $(A, i)$ be a CM simple abelian surface of type $(K, \Phi)$. Then there exists a unique locally constant homomorphism $u^{'}: \mathbb{A}^{\times}_{k}\to K^{\ast\times }$ satisfying the following properties:
          \\\indent (i) Fix an embedding $\tau: K^{\ast}\to \mathbb{C}$. The function
          \begin{align}
              \chi^{'}_{\tau}: \mathbb{A}^{\times}_{k}\xrightarrow{u^{'}(-)\cdot \overline{{\rm Nm}_{k/K^{\ast}}(-)}/{\rm Nm}_{k/K^{\ast}}(-)}\mathbb{A}^{\times}_{K^{\ast}} \xrightarrow{\tau-{\rm projection}}\mathbb{C}^{\times}\notag 
          \end{align}
          is an algebraic Hecke character.
          \\\indent (ii) Let 
          \begin{align}
              \rho_{l}: {\rm Gal}(\overline{k}/k)\to GL(T_{l}(A_{\overline{k}})) \notag 
          \end{align}
          be the $l$-adic representation. Then $\rho_{l}$ is diagonalized by the algebraic Hecke characters $\chi_{\tau}$, where $\tau$ runs over the set ${\rm Hom}(K^{\ast}, \mathbb{C})$.
      \end{cor}

      \begin{prop}\label{onegai} Let $(A, i)$ be a CM simple abelian surface of type $(K, \Phi)$. Let $X:={\rm Km}(A)$ be the Kummer surface, and let $\chi_{\tau}$ be the algebraic Hecke characters in Theorem \ref{main theorem} for $X$. 
          Then, for every embedding $\tau\in {\rm Hom}(K^{\ast}, \mathbb{C})$, we have  $\chi_{\tau}=\chi^{'}_{\tau}$. 
      \end{prop}

      \begin{proof}By the construction of Kummer surfaces, there is an isomorphism of ${\rm Gal}(\overline{k}/k)$-modules
      \begin{align}
          H^{2}(X_{\overline{k}}, \mathbb{Q}_{l})(1)\cong H^{2}(A_{\overline{k}}, \mathbb{Q}_{l})(1)\oplus V, \notag 
      \end{align}
      where $V$ is a 16-dimensional $\mathbb{Q}_{l}$-vector space and is generated by the cohomology classes of the exceptional divisors of the minimal resolution $X\to A/\langle-1\rangle$. In particular, this induces an isomorphism of ${\rm Gal}(\overline{k}/k)$-modules
          \begin{align}
              T(A_{\mathbb{C}})_{\mathbb{Q}_{l}}\cong T(X_{\mathbb{C}})_{\mathbb{Q}_{l}}.  \notag 
          \end{align}
          For any $\sigma\in{\rm Gal}(\overline{k}/k)$, let $s\in \mathbb{A}^{\times}_{k}$ be an id\'ele such that  $\sigma|_{k^{ab}}={\rm art}_{k}(s)$.  By Proposition \ref{hecke1}, the action of $\sigma$ on $T(X_{\mathbb{C}})_{\mathbb{Q}_{l}}$ is given by 
          \begin{align}u(t)\overline{{\rm Nm}_{k/K^{\ast}}(t)}/{\rm Nm}_{k/K^{\ast}}(t).\notag 
          \end{align}On the other hand, by  Proposition \ref{simple cm1}, the action of $\sigma$ on $T(A)_{\mathbb{Q}_{l}}$ is given by 
          \begin{align}
          u^{'}(t)\overline{{\rm Nm}_{k/K^{\ast}}(t)}/{\rm Nm}_{k/K^{\ast}}(t). \notag 
          \end{align}Hence we have
          \begin{align}
              \chi_{\tau}(t)=\left(u(t)\frac{\overline{{\rm Nm}_{k/K^{\ast}}(t)}}{{\rm Nm}_{k/K^{\ast}}(t)}\right)_{\tau}=\left(u^{'}(t)\frac{\overline{{\rm Nm}_{k/K^{\ast}}(t)}}{{\rm Nm}_{k/K^{\ast}}(t)}\right)_{\tau}=\chi^{'}_{\tau}(t).\notag 
          \end{align}
      \end{proof} 

\end{document}
